\documentclass[10pt,a4paper]{amsart}
\usepackage[margin=19mm]{geometry}
\usepackage{amsmath,amssymb,mathtools}
\usepackage[scr=rsfs,cal=euler]{mathalfa}
\usepackage{xfrac}
\usepackage[unicode,hidelinks,bookmarks=false]{hyperref}
\newcommand{\ie}{i.e.\ }
\newcommand{\cf}{cf.\ }
\newcommand{\resp}{respectively\ }
\newcommand{\Subsection}{\subsection}
\providecommand{\nobreakdash}{\nobreak\hbox{-}}
\setlength{\emergencystretch}{2em}
\multlinegap0pt
\usepackage{amssymb}
\newtheorem{theorem}{Theorem}
\title{Normal random variables}
\author{Teo Banica}
\date{Source: arXiv:2609.06859v1 (CC BY 4.0)}
\begin{document}
\maketitle

\noindent\emph{Source and licence.} Normal random variables. Version arXiv:2609.06859v1 (CC BY 4.0), distributed by arXiv under the Creative Commons Attribution 4.0 International licence, http://creativecommons.org/licenses/by/4.0/, as recorded in the arXiv licence metadata for this exact version and shown on the abstract page https://arxiv.org/abs/2609.06859v1. Full licence notice: \texttt{licenses/arxiv-2609.06859-CC-BY-4.0.txt}.\par\smallskip

\noindent\textit{Editorial scope.} Counted unit: the article's theorem theorem (source0.tex lines 9705--9709). The article's complete original proof of it is delivered with it (source0.tex lines 9711--9855, one \texttt{proof} environment of 145 lines). No mathematics is written, repaired or re-derived: the statement and the proof are the article's own lines, in the article's own notation, and no retained line is substituted or re-flowed.  No further source-local context is needed: every cross-reference of every retained line resolves inside the two delivered spans.  Nothing was omitted from any retained span, so the package carries no omission record.  No retained line contains a citation, so the wrapper carries no bibliography and no bibliography entry.  No retained line contains a figure environment, an image-inclusion command or drawing code, so no image is delivered and the package declares no asset dependency.  Editorial formatting only, in addition: an AMS \texttt{amsart} preamble that restates the article's own declarations and macros used by the retained lines, namely the environments \texttt{theorem} on separate counters, together with the standard packages those lines need.  The retained environments print in this document as Theorem 1 in order of appearance, whereas the article prints the same items with the numbering of its own full document.  The complete native source of the article as distributed in the arXiv e-print for this exact version is bundled unchanged as \texttt{sources/209580/source0.tex}.
\begin{theorem}[again]
The density of the chi squared law is given by
$$\chi_n^2=\frac{x^{n/2-1}e^{-x/2}}{2^{n/2}\Gamma(n/2)}\,dx$$
on $[0,\infty)$, this time coming via a direct density transport.
\end{theorem}

\begin{proof}
Buckle up, the computations will be a bit harder than expected:

\medskip

(1) Consider a vector $f=(f_i)$ whose components follow $g_1$ and are independent, so that $||f||^2\sim\chi_n^2$. We have the following formula, for any function $\varphi:[0,\infty)\to\mathbb R$:
$$E\left(\varphi(||f||^2)\right)
=\frac{1}{\sqrt{(2\pi)^n}}\int_{\mathbb R^n}\varphi(||x||^2)e^{-||x||^2/2}dx$$

Thus, we have our density, provided that we can compute the integral on the right.

\medskip

(2) In order to do so, we can use spherical coordinates. There are many possible conventions here, and we will use in this book the following nice-looking formulae:
$$\begin{cases}
x_1\!\!\!&=\ r\cos t_1\\
x_2\!\!\!&=\ r\sin t_1\cos t_2\\
\vdots\\
x_{n-1}\!\!\!&=\ r\sin t_1\sin t_2\ldots\sin t_{n-2}\cos t_{n-1}\\
x_n\!\!\!&=\ r\sin t_1\sin t_2\ldots\sin t_{n-2}\sin t_{n-1}
\end{cases}$$

We will also need the Jacobian of these coordinates. By developing over the last column, we obtain the following formula for this Jacobian $J_n$:
\begin{eqnarray*}
J_n
&=&r\sin t_1\ldots\sin t_{n-2}\sin t_{n-1}\times \sin t_{n-1}J_{n-1}\\
&+&r\sin t_1\ldots \sin t_{n-2}\cos t_{n-1}\times\cos t_{n-1}J_{n-1}\\
&=&r\sin t_1\ldots\sin t_{n-2}(\sin^2 t_{n-1}+\cos^2 t_{n-1})J_{n-1}\\
&=&r\sin t_1\ldots\sin t_{n-2}J_{n-1}
\end{eqnarray*}

Thus, by recurrence, we conclude that the Jacobian is given by:
$$J_n=r^{n-1}\sin^{n-2}t_1\sin^{n-3}t_2\,\ldots\,\sin^2t_{n-3}\sin t_{n-2}$$

Observe that this formula fits with the well-known formulae at $n=1,2,3$, that you surely know from physics class, namely $J_1=1$, $J_2=r$, $J_3=r^2\sin t_1$.

\medskip

(3) Getting back to our computation (1), this can be continued as follows, with the $2^n$ factor coming from the fact that we restricted attention to $x=(x_i)$ with $x_i\geq0$:
\begin{eqnarray*}
E\left(\varphi(||f||^2)\right)
&=&\frac{1}{\sqrt{(2\pi)^n}}\int_{\mathbb R^n}\varphi(||x||^2)e^{-||x||^2/2}dx\\
&=&\frac{2^n}{\sqrt{(2\pi)^n}}\int_0^{\pi/2}\ldots\int_0^{\pi/2}\int_0^\infty\varphi(r^2)e^{-r^2/2}\times r^{n-1}\\
&&\times\ \sin^{n-2}t_1\sin^{n-3}t_2\,\ldots\,\sin^2t_{n-3}\sin t_{n-2}\,drdt_1\ldots dt_{n-1}\\
&=&\left(\frac{2}{\pi}\right)^{n/2}\int_0^{\pi/2}\sin^{n-2}t_1dt_1\ldots\ldots\int_0^{\pi/2}\sin^0t_{n-1}dt_{n-1}\\
&&\times\int_0^\infty\varphi(r^2)r^{n-1}e^{-r^2/2}dr
\end{eqnarray*}

(4) Now if we denote by $C$ the above product of trigonometric integrals, with the normalization factor in front of it included, we have:
\begin{eqnarray*}
E\left(\varphi(||f||^2)\right)
&=&C\int_0^\infty\varphi(r^2)r^{n-1}e^{-r^2/2}dr\\
&=&C\int_0^\infty\varphi(y)y^{(n-1)/2}e^{-y/2}\frac{dy}{2\sqrt{y}}\\
&=&\frac{C}{2}\int_0^\infty\varphi(y)y^{n/2-1}e^{-y/2}dy
\end{eqnarray*}

(5) As a comment now, the trigonometric factor $C$ is more or less the same thing as the volume $V$ of the unit sphere, as shown by the following computation:
\begin{eqnarray*}
V
&=&2^n\int_0^{\pi/2}\ldots\int_0^{\pi/2}\int_0^1r^{n-1}\\
&&\times\ \sin^{n-2}t_1\sin^{n-3}t_2\,\ldots\,\sin^2t_{n-3}\sin t_{n-2}\,drdt_1\ldots dt_{n-1}\\
&=&\frac{2^n}{n}\int_0^{\pi/2}\sin^{n-2}t_1dt_1\ldots\ldots\int_0^{\pi/2}\sin^0t_{n-1}dt_{n-1}
\end{eqnarray*}

Even better, observe that the area $A$ of the unit sphere is given by:
$$A=nV=2^n\int_0^{\pi/2}\sin^{n-2}t_1dt_1\ldots\ldots\int_0^{\pi/2}\sin^0t_{n-1}dt_{n-1}$$

But with this, our computation in (3) simply reads, with $\psi(r)=\varphi(r^2)e^{-r^2/2}$:
$$\int_{\mathbb R^n}\psi(||x||)dx=A\int_0^\infty\psi(r)r^{n-1}dr$$

In fact, and we will leave this as an instructive geometry exercise, it is possible to come upon this latter formula, valid for any $\psi$, just by thinking. Enjoy.

\medskip

(6) Back to our business now, regardless of the method used, we need to compute that integrals of sines. I mean, even when short-circuiting the spherical coordinates, via some direct geometry, as suggested in (5), this would lead us into the computation of $V$. And here, regardless of the method used, which can be spherical coordinates, or recurrence by slicing, you are led in practice to computing the above integrals of sines. 

\medskip

(7) So, trigonometric integrals. Following Wallis, our claim is that we have the following formulae, where $\varepsilon(p)=1$ if $p$ is even, and $\varepsilon(p)=0$ if $p$ is odd:
$$\int_0^{\pi/2}\cos^pt\,dt=\int_0^{\pi/2}\sin^pt\,dt=\left(\frac{\pi}{2}\right)^{\varepsilon(p)}\frac{p!!}{(p+1)!!}$$

Indeed, our integrals being equal, let us look at the one on the left $I_p$. We have:
\begin{eqnarray*}
(\cos^pt\sin t)'
&=&p\cos^{p-1}t(-\sin t)\sin t+\cos^pt\cos t\\
&=&p\cos^{p+1}t-p\cos^{p-1}t+\cos^{p+1}t\\
&=&(p+1)\cos^{p+1}t-p\cos^{p-1}t
\end{eqnarray*}

By integrating now between $0$ and $\pi/2$, we obtain the following formula:
$$(p+1)I_{p+1}=pI_{p-1}$$

Thus we can compute $I_p$ by recurrence, and we obtain in this way:
\begin{eqnarray*}
I_p
&=&\frac{p-1}{p}\,I_{p-2}\\
&=&\frac{p-1}{p}\cdot\frac{p-3}{p-2}\,I_{p-4}\\
&\vdots&\\
&=&\frac{p!!}{(p+1)!!}\,I_{1-\varepsilon(p)}
\end{eqnarray*}

Now the initial data being $I_0=\frac{\pi}{2}$ and $I_1=1$, we obtain the result.

\medskip

(8) And with this, good news, we can finish. The factor $C$ from (4) is given by:
\begin{eqnarray*}
C
&=&\left(\frac{2}{\pi}\right)^{n/2}\int_0^{\pi/2}\sin^{n-2}t_1dt_1\ldots\ldots\int_0^{\pi/2}\sin^0t_{n-1}dt_{n-1}\\
&=&\left(\frac{2}{\pi}\right)^{n/2}\left(\frac{\pi}{2}\right)^{\varepsilon(n-2)+\ldots+\varepsilon(0)}\frac{(n-2)!!}{(n-1)!!}\cdot\frac{(n-3)!!}{(n-2)!!}\ldots\ldots\frac{0!!}{1!!}\\
&=&\left(\frac{2}{\pi}\right)^{n/2}\left(\frac{\pi}{2}\right)^{[n/2]}\frac{1}{(n-1)!!}\\
&=&\left(\frac{2}{\pi}\right)^{\varepsilon(n)/2}\frac{1}{(n-1)!!}
\end{eqnarray*}

(9) Thus, by getting back now to our formula in (4), that reads:
$$E\left(\varphi(||f||^2)\right)
=\frac{1}{2}\left(\frac{2}{\pi}\right)^{\varepsilon(n)/2}\frac{1}{(n-1)!!}\int_0^\infty\varphi(y)y^{n/2-1}e^{-y/2}dy$$

In order to further process the normalization factor, recall from the proof of Theorem 10.2 that we have the following formula, with $c=\sqrt{2},\sqrt{\pi}$ for $n$ even, odd:
$$\Gamma\left(\frac{n}{2}\right)=\frac{(n-1)!!}{2^{(n-1)/2}}\,c$$

Equivalently, with $\varepsilon(n)=1$ if $n$ is even, and $\varepsilon(n)=0$ if $n$ is odd, we have:
$$2^{n/2}\Gamma\left(\frac{n}{2}\right)=(n-1)!!\,\sqrt{2}\,c=(n-1)!!\,2\left(\frac{\pi}{2}\right)^{\varepsilon(n)/2}$$

Thus, in terms of the gamma function, our expectation formula reads:
$$E\left(\varphi(||f||^2)\right)=\frac{1}{2^{n/2}\Gamma(n/2)}\int_0^\infty\varphi(y)y^{n/2-1}e^{-y/2}dy$$

We are therefore led to the density in the statement.

\medskip

(10) Finally, as a byproduct of our study, let us record some useful formulae, for the volume $V$ and area $A$ of the unit sphere in $\mathbb R^n$, obtained via (5), thanks to the trigonometric integral computed in (8). According to our computations, these are:
$$V=\left(\frac{\pi}{2}\right)^{[n/2]}\frac{2^n}{(n+1)!!}\quad,\quad 
A=\left(\frac{\pi}{2}\right)^{[n/2]}\frac{2^n}{(n-1)!!}$$

Equivalently, in terms of the gamma function, the formulae are a bit more complicated, as follows, with our usual convention $c=\sqrt{2},\sqrt{\pi}$ for $n$ even, odd:
$$V=\left(\frac{\pi}{2}\right)^{[n/2]}\frac{2^{(n-1)/2}c}{\Gamma(n/2+1)}\quad,\quad 
A=\left(\frac{\pi}{2}\right)^{[n/2]}\frac{2^{(n+1)/2}c}{\Gamma(n/2)}$$

Finally, for the story to be complete, let us record as well the Stirling asymptotics:
$$V\simeq\left(\frac{2\pi e}{n}\right)^{n/2}\frac{1}{\sqrt{\pi n}}
\quad,\quad A\simeq\left(\frac{2\pi e}{n}\right)^{n/2}\sqrt{\frac{n}{\pi}}$$

And with of course exercise for you, to learn more about all this, spheres.
\end{proof}

\end{document}
